\chapter*{Lampiran: Ringkasan Rumus Matematika dan Konteks Data}
\addcontentsline{toc}{chapter}{Lampiran: Ringkasan Rumus Matematika}

Dokumen ini merangkum notasi matematika dan algoritma STATIS (\textit{Structuration des Tableaux à Trois Indices de la Statistique}) serta \textit{Dynamic Biplot} yang dikaitkan secara langsung dengan variabel data penelitian (Kabupaten/Kota di Banten, DKI Jakarta, dan Jawa Barat tahun 2020--2024).

\section*{1. Definisi Data Tiga-Arah (Data Panel)}

Data penelitian disusun dalam struktur tensor atau larik tiga dimensi $\mathcal{X}$.

\begin{itemize}
    \item \textbf{Unit Observasi ($n=41$):} Merepresentasikan 41 Kabupaten/Kota administratif.
    \item \textbf{Variabel ($p=12$):} Merepresentasikan 12 indikator sosio-ekonomi (misalnya: IPM, PDRB per kapita, Tingkat Pengangguran, Angka Harapan Hidup, dll).
    \item \textbf{Dimensi Waktu ($K=5$):} Merepresentasikan periode tahun 2020, 2021, 2022, 2023, dan 2024.
\end{itemize}

\textbf{Notasi Data:}
\begin{equation}
    x_{ijt}
\end{equation}
\textit{Interpretasi:} Nilai indikator ke-$j$ (misal: IPM) untuk Kabupaten/Kota ke-$i$ (misal: Kota Bandung) pada tahun ke-$t$ (misal: 2020).

\textbf{Matriks Tahun Tertentu:}
\begin{equation}
    \mathbf{X}_t \quad (t = 2020, \dots, 2024)
\end{equation}
\textit{Interpretasi:} Sebuah matriks berukuran $41 \times 12$ yang berisi potret lengkap kondisi sosio-ekonomi seluruh wilayah pada satu tahun spesifik.

\section*{2. Standardisasi Data (Z-Score)}

Karena indikator memiliki satuan yang sangat beragam (contoh: PDRB dalam Rupiah jutaan vs Tingkat Kemiskinan dalam Persen), dilakukan penyetaraan skala.

\begin{equation}
    z_{ijt} = \frac{x_{ijt} - \bar{x}_{jt}}{s_{jt}}
\end{equation}

\begin{itemize}
    \item $\bar{x}_{jt}$: Rata-rata indikator $j$ pada tahun $t$ (Rata-rata IPM seluruh kota tahun 2020).
    \item $s_{jt}$: Standar deviasi indikator $j$ pada tahun $t$.
    \item \textit{Implementasi Python:} \texttt{final.py} Baris 316 (\texttt{scaler.fit\_transform}).
    \item \textit{Tujuan:} Agar fluktuasi PDRB (jutaan) tidak mendominasi fluktuasi Kemiskinan (persen) dalam perhitungan jarak.
\end{itemize}

\section*{3. Analisis STATIS (Stabilitas Struktur Ekonomi)}

\subsection*{Tahap 1: Intra-struktur (Peta Kemiripan Tahunan)}
Mengubah data angka menjadi konfigurasi jarak antar-kota.

\begin{equation}
    \mathbf{W}_t = \mathbf{Z}_t \mathbf{Z}_t^T \quad (t=2020, \dots, 2024)
\end{equation}
\textit{Implementasi Python:} \texttt{final.py} Baris 319 \\ (\texttt{Ws.append(arr\_scaled @ arr\_scaled.T)}).

\textit{Interpretasi:} Matriks $\mathbf{W}_{2020}$ ($41 \times 41$) menggambarkan peta persaingan/kemiripan ekonomi antar distrik pada tahun awal \\ pandemi. Jika Kota A dan Kota B memiliki profil ekonomi mirip di 2020, maka nilai elemen $(A,B)$ pada matriks ini akan tinggi.

% Wraps generated files
% Generated by final.py
\newcommand{\statisDatasetName}{Data Statistik}
\newcommand{\statisAlphaTahunA}{0.2006} % 2020
\newcommand{\statisAlphaTahunB}{0.1993} % 2021
\newcommand{\statisAlphaTahunC}{0.1985} % 2022
\newcommand{\statisAlphaTahunD}{0.2017} % 2023
\newcommand{\statisAlphaTahunE}{0.1999} % 2024
\newcommand{\statisEigenA}{238.8721}
\newcommand{\statisInertiaA}{48.55}
% Generated by final.py (Permutation)
\newcommand{\statisPValue}{0.0010}
\newcommand{\statisKet}{Signifikan (Struktur Stabil)}


\begin{table}[H]
\centering
\caption{Nilai Bobot (Alpha)} \label{tab:alpha}
\begin{tabular}{|c|c|}
\hline
\textbf{Tahun} & \textbf{Bobot (Alpha)} \\
\hline
2020 & 0.2006 \\
2021 & 0.1993 \\
2022 & 0.1985 \\
2023 & 0.2017 \\
2024 & 0.1999 \\
\hline
\end{tabular}
\end{table}


\subsection*{Tahap 2: Inter-struktur (Hubungan Antar-Tahun)}
Mengukur seberapa mirip "peta ekonomi" tahun 2020 dibandingkan 2021, 2022, dst.

\textbf{Koefisien RV (Korelasi Waktu):}
\begin{equation}
    RV(\mathbf{W}_{2020}, \mathbf{W}_{2021}) = \frac{tr(\mathbf{W}_{2020} \mathbf{W}_{2021})}{\sqrt{tr(\mathbf{W}_{2020}^2) tr(\mathbf{W}_{2021}^2)}}
\end{equation}
\textit{Implementasi Python:} \texttt{final.py} Baris 109 (\texttt{num = np.trace(Ws\_list[i].T @ Ws\_list[j])}).
\textit{Interpretasi:} Jika $RV \approx 1$, artinya struktur ketimpangan antar-wilayah tidak berubah dari 2020 ke 2021.

\begin{table}[H]
\centering
\caption{Koefisien RV Antar-tahun} \label{tab:rv}
\begin{tabular}{|c|c|c|c|c|c|}
\hline
\textbf{2020} & \textbf{2021} & \textbf{2022} & \textbf{2023} & \textbf{2024} & \textbf{Tahun} \\
\hline
1.0000 & 0.9533 & 0.9238 & 0.9583 & 0.9471 & 2020 \\
0.9533 & 1.0000 & 0.9382 & 0.9366 & 0.9245 & 2021 \\
0.9238 & 0.9382 & 1.0000 & 0.9456 & 0.9255 & 2022 \\
0.9583 & 0.9366 & 0.9456 & 1.0000 & 0.9671 & 2023 \\
0.9471 & 0.9245 & 0.9255 & 0.9671 & 1.0000 & 2024 \\
\hline
\end{tabular}
\end{table}


\textbf{Bobot Tahun ($\alpha_t$):}
Dihitung dari dekomposisi matriks korelasi antar-tahun.
\begin{equation}
    \mathbf{W}_c = \alpha_{2020}\mathbf{W}_{2020} + \dots + \alpha_{2024}\mathbf{W}_{2024}
\end{equation}
\textit{Implementasi Python:} \texttt{final.py} Baris 337 (\texttt{Wc += a * W}).

\textit{Interpretasi:} $\alpha_t$ menentukan seberapa "penting" suatu tahun. Jika tahun 2020 (awal pandemi) memiliki pola yang sangat aneh/berbeda \\ (anomali) dibanding 4 tahun lainnya, algoritma akan memberikan bobot $\alpha_{2020}$ yang kecil secara otomatis agar tidak merusak representasi rata-rata (kompromi).
Berikut adalah nilai bobot yang dihasilkan dari pengolahan data \textbf{\statisDatasetName}:
\begin{align*}
    \alpha_{2020} &= \statisAlphaTahunA \\
    \alpha_{2021} &= \statisAlphaTahunB \\
    \alpha_{2022} &= \statisAlphaTahunC \\
    \alpha_{2023} &= \statisAlphaTahunD \\
    \alpha_{2024} &= \statisAlphaTahunE
\end{align*}
\textit{Interpretasi:} Nilai bobot ini menunjukkan kontribusi setiap tahun terhadap pembentukan struktur kompromi.

\subsection*{Tahap 3: Struktur Kompromi (Peta Ekonomi Rata-rata)}
\begin{equation}
    \mathbf{F}_c = \text{EigDecomp}(\mathbf{W}_c)
\end{equation}
\textit{Interpretasi:} Koordinat $\mathbf{F}_c$ adalah posisi rata-rata setiap Kabupaten/Kota selama 5 tahun. Ini menjadi peta dasar (fonds de carte) \\ untuk melihat posisi relatif wilayah secara umum (misal: Kelompok Jabodetabek vs Non-Jabodetabek).
Dari hasil analisis, dimensi pertama (Sumbu X) mampu menjelaskan variansi sebesar \textbf{\statisInertiaA\%} dengan nilai eigen \textbf{\statisEigenA}.

\begin{table}[H]
\centering
\caption{Dekomposisi Nilai Eigen} \label{tab:eigen}
\begin{tabular}{|c|c|c|c|c|}
\hline
\textbf{Dim} & \textbf{Eigen} & \textbf{Akar Ciri} & \textbf{Var (\%)} & \textbf{Kum (\%)} \\
\hline
Dim 1 & 238.8721 & 15.4555 & 48.5500 & 48.5500 \\
Dim 2 & 79.0492 & 8.8910 & 16.0700 & 64.6200 \\
Dim 3 & 58.3429 & 7.6382 & 11.8600 & 76.4800 \\
Dim 4 & 21.6880 & 4.6570 & 4.4100 & 80.8800 \\
Dim 5 & 17.6519 & 4.2014 & 3.5900 & 84.4700 \\
Dim 6 & 16.4978 & 4.0617 & 3.3500 & 87.8300 \\
Dim 7 & 15.4929 & 3.9361 & 3.1500 & 90.9700 \\
Dim 8 & 11.1258 & 3.3355 & 2.2600 & 93.2400 \\
Dim 9 & 7.7464 & 2.7832 & 1.5700 & 94.8100 \\
Dim 10 & 6.0717 & 2.4641 & 1.2300 & 96.0400 \\
Dim 11 & 4.2161 & 2.0533 & 0.8600 & 96.9000 \\
Dim 12 & 2.8697 & 1.6940 & 0.5800 & 97.4800 \\
Dim 13 & 2.3563 & 1.5350 & 0.4800 & 97.9600 \\
Dim 14 & 2.2951 & 1.5150 & 0.4700 & 98.4300 \\
Dim 15 & 1.6780 & 1.2954 & 0.3400 & 98.7700 \\
Dim 16 & 1.4369 & 1.1987 & 0.2900 & 99.0600 \\
Dim 17 & 0.9965 & 0.9983 & 0.2000 & 99.2700 \\
Dim 18 & 0.7466 & 0.8641 & 0.1500 & 99.4200 \\
Dim 19 & 0.5621 & 0.7497 & 0.1100 & 99.5300 \\
Dim 20 & 0.5185 & 0.7201 & 0.1100 & 99.6400 \\
Dim 21 & 0.4585 & 0.6771 & 0.0900 & 99.7300 \\
Dim 22 & 0.4070 & 0.6380 & 0.0800 & 99.8100 \\
Dim 23 & 0.3142 & 0.5605 & 0.0600 & 99.8800 \\
Dim 24 & 0.1573 & 0.3966 & 0.0300 & 99.9100 \\
Dim 25 & 0.1315 & 0.3626 & 0.0300 & 99.9400 \\
Dim 26 & 0.0922 & 0.3037 & 0.0200 & 99.9500 \\
Dim 27 & 0.0678 & 0.2604 & 0.0100 & 99.9700 \\
Dim 28 & 0.0470 & 0.2169 & 0.0100 & 99.9800 \\
Dim 29 & 0.0332 & 0.1822 & 0.0100 & 99.9800 \\
Dim 30 & 0.0190 & 0.1378 & 0.0000 & 99.9900 \\
Dim 31 & 0.0149 & 0.1221 & 0.0000 & 99.9900 \\
Dim 32 & 0.0138 & 0.1176 & 0.0000 & 99.9900 \\
Dim 33 & 0.0068 & 0.0823 & 0.0000 & 100.0000 \\
Dim 34 & 0.0060 & 0.0775 & 0.0000 & 100.0000 \\
Dim 35 & 0.0048 & 0.0696 & 0.0000 & 100.0000 \\
Dim 36 & 0.0036 & 0.0602 & 0.0000 & 100.0000 \\
Dim 37 & 0.0034 & 0.0580 & 0.0000 & 100.0000 \\
Dim 38 & 0.0021 & 0.0458 & 0.0000 & 100.0000 \\
Dim 39 & 0.0013 & 0.0364 & 0.0000 & 100.0000 \\
Dim 40 & 0.0012 & 0.0352 & 0.0000 & 100.0000 \\
Dim 41 & 0.0000 & 0.0000 & 0.0000 & 100.0000 \\
\hline
\end{tabular}
\end{table}
\begin{longtable}{|c|c|c|}
\caption{Koordinat Skor Kompromi} \label{tab:scores} \\
\hline
\textbf{Kab/Kota} & \textbf{PC1} & \textbf{PC2} \\
\hline
\endfirsthead
\hline
\textbf{Kab/Kota} & \textbf{PC1} & \textbf{PC2} \\
\hline
\endhead
\hline
\endfoot
Kep. Seribu & 0.8923 & -1.2073 \\
Kota Jakarta Selatan & -4.8731 & -0.4667 \\
Kota Jakarta Timur & -3.7890 & 0.0126 \\
Kota Jakarta Pusat & -3.9814 & -0.0467 \\
Kota Jakarta Barat & -3.3486 & -0.3053 \\
Kota Jakarta Utara & -2.9407 & 0.3065 \\
Bogor & 0.4026 & 0.0378 \\
Sukabumi & 1.9922 & 0.3459 \\
Cianjur & 3.0050 & 0.6465 \\
Bandung & -0.2850 & 1.3296 \\
Garut & 2.5490 & 0.7424 \\
Tasikmalaya & 3.0968 & 1.1724 \\
Ciamis & 1.7104 & 1.2463 \\
Kuningan & 1.2773 & 1.3932 \\
Cirebon & 2.1439 & 1.3411 \\
Majalengka & 2.5288 & 1.0199 \\
Sumedang & 0.5391 & 1.1202 \\
Indramayu & 2.1079 & 0.8393 \\
Subang & 1.5388 & 1.1650 \\
Purwakarta & 1.0637 & 0.2658 \\
Karawang & 0.2501 & 0.2822 \\
Bekasi & -1.4773 & 0.0343 \\
Bandung Barat & 1.9876 & 1.4169 \\
Pangandaran & 2.2717 & 1.6373 \\
Kota Bogor & -1.9711 & 0.8351 \\
Kota Sukabumi & -1.4078 & 0.4201 \\
Kota Bandung & -3.0860 & 0.6713 \\
Kota Cirebon & -0.7971 & 1.3383 \\
Kota Bekasi & -3.6009 & 0.2417 \\
Kota Depok & -3.5738 & -0.3413 \\
Kota Cimahi & -2.2576 & 0.2520 \\
Kota Tasikmalaya & 0.3450 & 1.2384 \\
Kota Banjar & 1.3293 & 0.2293 \\
Pandeglang & 3.7767 & -3.2485 \\
Lebak & 4.0257 & -1.5961 \\
Tangerang & 0.3939 & -1.8870 \\
Serang & 2.4297 & -4.1130 \\
Kota Tangerang & -2.0499 & -1.4172 \\
Kota Cilegon & 0.2294 & -3.2300 \\
Kota Serang & 0.6823 & -2.3970 \\
Kota Tangerang Selatan & -3.1301 & -1.3255 \\
\hline
\end{longtable}


\section*{4. Dynamic Biplot (Visualisasi Perubahan)}

Melihat bagaimana posisi Kabupaten/Kota "bergerak" dari tahun ke tahun di atas peta kompromi.

\textbf{Proyeksi Trajektori:}
\begin{equation}
    \mathbf{F}_{i,t} = \mathbf{Z}_{i,t} \mathbf{V}
\end{equation}
\textit{Implementasi Python:} \texttt{final.py} Baris 372 (\texttt{per\_year\_coords[year] = dfstd.values @ loadings}).

\begin{table}[H]
\centering
\caption{Loadings Kompromi} \label{tab:loadings}
\begin{tabular}{|c|c|c|}
\hline
\textbf{Variabel} & \textbf{PC1} & \textbf{PC2} \\
\hline
AHH\_L & -0.6913 & 0.6786 \\
AHH\_P & -0.6809 & 0.6912 \\
GK & -0.8438 & -0.2275 \\
GR & -0.7011 & 0.4376 \\
IPG & -0.7780 & -0.1173 \\
IPM & -0.9860 & -0.0235 \\
PDRB\_ADHK & -0.6468 & -0.0072 \\
PKKP & 0.5904 & 0.7154 \\
PPM & 0.7213 & 0.3949 \\
RLS & -0.9540 & -0.0926 \\
TPAK & 0.4467 & 0.3159 \\
TPT & -0.2072 & -0.2391 \\
\hline
\end{tabular}
\end{table}


\section*{5. Uji Permutasi (Validasi Statistik)}
Memastikan bahwa pola kestabilan yang ditemukan bukan kebetulan semata.
\begin{itemize}
    \item $H_0$: Struktur ekonomi antar-tahun bersifat acak (tidak ada pola tetap).
    \item \textbf{Statistik Uji ($p$-value):}
    \begin{equation}
        p = \frac{\sum_{i=1}^{N_{perm}} I(\lambda^{(perm)}_i \ge \lambda^{(obs)}) + 1}{N_{perm} + 1}
    \end{equation}
    \textit{Implementasi Python:} \texttt{final.py} Baris 453.
    
    \item Jika $p < 0.05$: Terdapat struktur ekonomi yang stabil dan signifikan dari 2020 hingga 2024 di wilayah studi (Banten, DKI, Jabar).
    \item \textbf{Hasil Uji:} Diperoleh $p$-value = \textbf{\statisPValue}. 
    \item \textbf{Kesimpulan:} \statisKet.
\end{itemize}

\section*{6. Uji Autokorelasi Spasial (Global Moran's~$I$)}
Mengukur apakah wilayah yang berdekatan memiliki karakteristik ekonomi yang mirip (pengelompokan spasial) pada Skor Kompromi.

\begin{equation}
    I = \frac{N}{\sum_i \sum_j w_{ij}} \frac{\sum_i \sum_j w_{ij} (x_i - \bar{x})(x_j - \bar{x})}{\sum_i (x_i - \bar{x})^2}
\end{equation}

\begin{itemize}
    \item $N$: Jumlah wilayah (41).
    \item $x_i$: Nilai Skor Kompromi (PC1) wilayah $i$.
    \item $w_{ij}$: Elemen matriks pembobot spasial (bernilai 1 jika wilayah $i$ dan $j$ bertetangga/bersentuhan, 0 jika tidak). \textbf{Catatan:} Matriks $w_{ij}$ ini berbeda dengan matriks $\mathbf{W}_t$ (produk skalar) pada analisis STATIS.
    \item \textit{Implementasi Python:} \texttt{final.py} Baris 119 (\texttt{\_calculate\_moran\_manual}).
    \item \textit{Interpretasi:} Jika $I > E[I]$, pola cenderung berkelompok (\textit{Clustered}).
\end{itemize}

\section*{7. Clustering Trajektori (K-Means)}
Mengelompokkan Kabupaten/Kota berdasarkan kemiripan pola pergerakan selama 5 tahun.

\textbf{Fungsi Objektif K-Means:}
Meminimalkan variansi dalam cluster (Within-Cluster Sum of Squares).
\begin{equation}
    J = \sum_{j=1}^{k} \sum_{i \in C_j} ||\mathbf{x}_i^{(traj)} - \boldsymbol{\mu}_j||^2
\end{equation}
\textit{Implementasi Python:} \texttt{final.py} Baris 1146 (\texttt{KMeans(n\_clusters=k\_clusters)}).

\textbf{Hierarchical Clustering (Ward):}
Digunakan untuk dendrogram, meminimalkan peningkatan variansi saat penggabungan cluster.
\begin{equation}
    d(u,v) = \sqrt{\frac{|u| \cdot |v|}{|u| + |v|}} ||\boldsymbol{\mu}_u - \boldsymbol{\mu}_v||^2
\end{equation}
\textit{Implementasi Python:} \texttt{final.py} Baris 1153 (\texttt{linkage(X\_traj, method='ward')}).
